
        \documentclass[fleqn]{article}      
        \usepackage{amsmath}
        \usepackage{fontspec}
        \usepackage{kotex} % 한국어 지원  

        \begin{document}      
         객관식 질문

1. 다음 중 적분의 정의와 가장 관련이 깊은 것은 무엇인가요?
   a) 미분  
   b) 극한  
   c) 면적  
   d) 함수  

2. 다음 중 부정적분의 표기가 올바른 것은 무엇인가요?
   a) ∫f(x)dx  
   b) ∫f(x)dx + C  
   c) ∫f(x)  
   d) ∫f(x) + C  

3. ∫x dx의 결과는 무엇인가요?
   a) x^2/2 + C  
   b) x^2 + C  
   c) 2x + C  
   d) x + C  

4. ∫(3x^2)dx의 결과는 무엇인가요?
   a) x^3 + C  
   b) x^3/3 + C  
   c) x^3/2 + C  
   d) 3x^3 + C  

5. 다음 중 정적분의 정의에 해당하는 것은 무엇인가요?
   a) ∫[a, b] f(x)dx  
   b) ∫f(x)dx  
   c) ∫f'(x)dx  
   d) ∫f(x)dx + C  

 주관식 질문

1. 적분의 기본 정리를 설명하세요.
2. 정적분과 부정적분의 차이를 설명하세요.
3. ∫x^2 dx의 결과를 구해보세요.
4. 정적분을 이용하여 구간 [1, 3]에서 f(x) = 2x의 면적을 구하세요.
5. ∫(2x + 1)dx의 결과를 구해보세요.
6. 평균값 정리의 개념을 설명하세요.
7. 적분을 활용하여 곡선 아래의 면적을 구하는 방법을 설명하세요.
8. ∫(x^3 - 4x + 1)dx의 결과를 구하세요.
9. 정적분의 기하학적 의미를 설명하세요.
10. 두 함수 f(x)와 g(x)의 교차점에서의 적분을 어떻게 구하는지 설명하세요.

 쉬운 문제

1. ∫(2x)dx를 구하세요.
2. ∫(5)dx의 결과는 무엇인가요?
3. ∫(x)dx 결과를 쓰세요.
4. ∫(3x + 2)dx의 결과를 구하세요.
5. ∫(x^2 + 1)dx의 결과를 구하세요.

 중간 문제

1. ∫(4x^3 - 2x^2 + x)dx를 구하세요.
2. f(x) = x^2의 정적분을 [0, 2] 구간에서 계산하세요.
3. ∫(sin x)dx의 결과를 구하세요.
4. ∫(e^x)dx의 결과는 무엇인가요?
5. ∫(1/x)dx의 결과를 구하세요.

 어려운 문제

1. ∫(x^4 - 3x^2 + 2)dx를 구하세요.
2. f(x) = x^3의 정적분을 [1, 4] 구간에서 계산하세요.
3. ∫(cos^2 x)dx의 결과를 구하세요.
4. ∫(x^2 * e^x)dx를 부분적분을 사용하여 구하세요.
5. ∫(tan x)dx의 결과를 구하세요.

      
        \end{document} 
    